\section{Governor: Stable Nonzero Scalar Occupancy}
\label{sec:governor}

Assume a scalar whose gauge-invariant radial variable is $s=|\phi|^2\geq0$, and restrict the leading potential to $V(s)=as+bs^2$ with $V(0)=0$. Phase invariance permits functions of $s$; it does not by itself select a polynomial of degree two. The dimension-four restriction defines the leading ansatz rather than an exact consequence of finite state count.

\subsection{The Sign Selection Theorem}
\nolinkurl{Governor.governor_selects_hat} proves an equivalence. Some positive occupancy costs less than the empty state, and the potential attains a minimum at positive occupancy, precisely when $a<0<b$. Writing $a=-\mu^2$, $b=\lambda$ gives
\begin{equation}
 V(\phi)=-\mu^2|\phi|^2+\lambda|\phi|^4,
 \qquad\mu^2>0,\quad\lambda>0.
 \label{eq:higgs_potential}
\end{equation}
The no-starvation and no-runaway conditions are explicit selection hypotheses. This theorem fixes the signs within the ansatz, not the numerical coefficients or the physical identity of the scalar.

\subsection{The Minimum and Field Convention}
Completing the square gives
\begin{equation}
 V(s)+\frac{\mu^4}{4\lambda}
 =\lambda\left(s-\frac{\mu^2}{2\lambda}\right)^2.
\end{equation}
Lean proves the lower bound and its equality condition (\nolinkurl{Potential}). With the convention $\langle|\phi|^2\rangle=v^2/2$, the minimum condition is $\mu^2=\lambda v^2$. This avoids confusing $v$ with the radial amplitude $v/\sqrt2$.

The candidate kinetic term is $-|D_\mu\phi|^2$ in the stated metric convention. Its coefficient, representation and covariant derivative must be derived from the same microscopic response used for the fermion and gauge sectors. Canonical normalization and a physical scalar mass do not follow from the potential minimum alone. Paper III tests proposed coefficients, including their joint mass relations; they are not assumed uniquely fixed here.

\subsection{Finite Spacing and the Ultraviolet Question}
\label{subsec:uv_completion}
Finite lattice resolution can regulate a specified dynamics, but a regulator alone does not prove naturalness, the observed hierarchy, or exact absence of higher operators. The Lean finite-amplitude counterexample retains an independent sixth-order interaction. Higher-dimensional operators with coefficients $c\ell^k$ vanish in the stated scaling limit; their finite-spacing effects require bounds.

Deriving the scalar kinetic response, its allowed higher terms and its relation to the detuning spectrum is required before claiming a physical Higgs sector or a resolution of the hierarchy problem. The Planck/Higgs ratio is not proved in this paper.


\subsection{The Link Response and Its Remaining Normalization}
The same specified comparison functional extends to transported amplitudes as
$C_{\rm link}=(c_0m/4)\lVert\phi'-U\phi\rVert^2$
(\nolinkurl{CommonResponse}). Sharing the prefactor is conditional on sharing the
readout and primitive repair multiplicity; this algebra does not prove that every
physical sector realizes the same repair process. Unnormalized field amplitudes
also require a map to maintained workload, rather than treating an arbitrary
squared norm as a probability for one event.

For a Hermitian generator in a unitary eigenbasis,
$\Theta=V\operatorname{diag}(\lambda)V^\dagger$, let
$U_\epsilon=V\operatorname{diag}(e^{i\epsilon\lambda})V^\dagger$ and
$c=V^\dagger\phi$. When $|\epsilon\lambda_k|\leq1$, \nolinkurl{LinkExpansion}
proves the controlled bound
\begin{equation}
 \left|\lVert\phi-U_\epsilon\phi\rVert^2-
 \epsilon^2\lVert\Theta\phi\rVert^2\right|
 \leq\frac5{48}\epsilon^4\sum_k |c_k|^2\lambda_k^4.
\end{equation}
Thus the leading link penalty used for the scalar/gauge mass response has an actual
finite-angle error bound. Selecting its physical representation, scalar kinetic
normalization, field-to-workload map and fermion coupling still requires the common
microscopic carrier dynamics.
